\subsection{\problemblock{*Dwn/crs} Data Block}\label{sec:block-dwn-crs}

East, north, downrange, and crossrange are defined from a reference point given by
longitude and geodetic latitude, as shown in
\cref{fig:downrange-crossrange-definition}. The trajectory is projected onto the
earth's surface, and the distances are measured along that surface curve as described
in \cref{sec:ranges-distances}. East is the component of distance traveled along a
constant latitude, and north is the component traveled along a constant longitude.
Downrange is the component traveled along a reference azimuth or heading, and
crossrange is the component traveled perpendicular to that heading.

\begin{figure}[htbp]
  \centering
  \begin{tikzpicture}[x=1cm,y=1cm,>=Latex,font=\small]
  \draw[->] (-0.3,0) -- (8,0) node[right] {longitude};
  \draw[->] (0,-0.3) -- (0,4.3) node[above] {latitude};
  \coordinate (R) at (1,0.8);
  \coordinate (P) at (5.4,2.5);
  \coordinate (V) at (6.7,3.35);
  \fill (R) circle (1.5pt) node[below left] {Reference point};
  \fill (V) circle (1.5pt) node[above right] {Vehicle};
  \draw[thick,->] (R) -- ($(R)!1.35!(P)$) node[pos=0.55,below] {downrange};
  \draw[thick] (P) -- (V) node[midway,above left] {crossrange};
  \draw[dashed] (R) -- ++(0,2.0);
  \draw[->] ($(R)+(0,0.65)$) arc[start angle=90,end angle=24,radius=0.65];
  \node at (1.55,1.42) {reference azimuth};
  \draw ($(P)+(-0.18,0.1)$) -- ++(0.17,0.3) -- ++(0.3,-0.17);
\end{tikzpicture}

  \caption{Downrange and crossrange definition.}
  \label{fig:downrange-crossrange-definition}
\end{figure}

All of these distances use the same reference point. The reference heading is used
only by the downrange and crossrange calculations. The reference longitude,
latitude, and heading default to the beginning of the trajectory given by the initial
conditions.

If a different reference point is required, or if multiple trajectories need to use the
same reference point, the values can be given in a \problemblock{*dwn/crs} data
block. For example,

\begin{lstlisting}[style=taosmanual]
*dwn/crs latgd=10 long=95 azm=45
\end{lstlisting}

places the reference point at a latitude of $10^\circ$~N and a longitude of
$95^\circ$~E. The reference heading or azimuth is $45^\circ$ from north, or
northeast.
